\documentclass[a4paper]{article}     
\usepackage{amsfonts}
\usepackage{amsmath}
\usepackage{amssymb}
\usepackage{graphicx}

\begin{document}
	
	\noindent \large Auteur: Victor Pena \\
	\noindent \large Studierichting: Informatica\\
	\noindent \large Oefeningenreeks 3A nr. 26.2\\
	
	\medskip
	
	\normalsize 
	
	$\operatorname{Bgsin}(2x^2+1)$ is alleen gedefinie{\"e}rd in een enkel punt $x=0$. De afgeleide bestaat niet, maar als we toch proberen die te vinden:\\
	
	$\dfrac{d}{dx}\operatorname{Bgsin}\left(2x^2+1\right) = \dfrac{1}{\sqrt{1-(2x^2+1)^2}}\cdot\dfrac{d}{dx}\left(2x^2+1\right)=\dfrac{4x 	}{\sqrt{1-(2x^2+1)^2}} =$ \\
	
	$\dfrac{4x}{\sqrt{-4x^4-4x^2}}=\dfrac{4x}{\sqrt{4x^2(-x^2-1})}=\dfrac{4x}{\sqrt{4x^2}\sqrt{-x^2-1}}=\dfrac{4x}{2x\sqrt{-x^2-1}}$ \\
	
	$ = \dfrac{2}{\sqrt{-x^2-1}} $ en die bestaat helemaal nergens!
	
	
	\begin{thebibliography}{999}
		%\bibitem{a} Referentie
	\end{thebibliography}
\end{document} 
